\documentclass[10pt,a4paper]{amsart}
\usepackage{amsmath,amssymb,amsthm,graphicx}
\usepackage{verbatim}
\newenvironment{reference}{\comment}{\endcomment}
%\newenvironment{reference}{}{}
\newenvironment{slogan}{\comment}{\endcomment}
\newenvironment{history}{\comment}{\endcomment}
% For commutative diagrams we use Xy-pic
\usepackage[all]{xy}

% We use 2cell for 2-commutative diagrams.
\xyoption{2cell}
% Theorem environments.
%
\theoremstyle{plain}
\newtheorem{theorem}[subsection]{Theorem}
\newtheorem{proposition}[subsection]{Proposition}
\newtheorem{lemma}[subsection]{Lemma}

\theoremstyle{definition}
\newtheorem{definition}[subsection]{Definition}
\newtheorem{example}[subsection]{Example}
\newtheorem{exercise}[subsection]{Exercise}
\newtheorem{situation}[subsection]{Situation}

\theoremstyle{remark}
\newtheorem{remark}[subsection]{Remark}
\newtheorem{remarks}[subsection]{Remarks}

\numberwithin{equation}{subsection}

% Macros
%
\def\lim{\mathop{\mathrm{lim}}\nolimits}
\def\colim{\mathop{\mathrm{colim}}\nolimits}
\def\Spec{\mathop{\mathrm{Spec}}}
\def\Hom{\mathop{\mathrm{Hom}}\nolimits}
\def\Ext{\mathop{\mathrm{Ext}}\nolimits}
\def\SheafHom{\mathop{\mathcal{H}\!\mathit{om}}\nolimits}
\def\SheafExt{\mathop{\mathcal{E}\!\mathit{xt}}\nolimits}
\def\Sch{\mathit{Sch}}
\def\Mor{\mathop{\mathrm{Mor}}\nolimits}
\def\Ob{\mathop{\mathrm{Ob}}\nolimits}
\def\Sh{\mathop{\mathit{Sh}}\nolimits}
\def\NL{\mathop{N\!L}\nolimits}
\def\CH{\mathop{\mathrm{CH}}\nolimits}
\def\proetale{{pro\text{-}\acute{e}tale}}
\def\etale{{\acute{e}tale}}
\def\QCoh{\mathit{QCoh}}
\def\Ker{\mathop{\mathrm{Ker}}}
\def\Im{\mathop{\mathrm{Im}}}
\def\Coker{\mathop{\mathrm{Coker}}}
\def\Coim{\mathop{\mathrm{Coim}}}

% Boxtimes
%
\DeclareMathSymbol{\boxtimes}{\mathbin}{AMSa}{"02}

%
% Macros for moduli stacks/spaces
%
\def\QCohstack{\mathcal{QC}\!\mathit{oh}}
\def\Cohstack{\mathcal{C}\!\mathit{oh}}
\def\Spacesstack{\mathcal{S}\!\mathit{paces}}
\def\Quotfunctor{\mathrm{Quot}}
\def\Hilbfunctor{\mathrm{Hilb}}
\def\Curvesstack{\mathcal{C}\!\mathit{urves}}
\def\Polarizedstack{\mathcal{P}\!\mathit{olarized}}
\def\Complexesstack{\mathcal{C}\!\mathit{omplexes}}
% \Pic is the operator that assigns to X its picard group, usage \Pic(X)
% \Picardstack_{X/B} denotes the Picard stack of X over B
% \Picardfunctor_{X/B} denotes the Picard functor of X over B
\def\Pic{\mathop{\mathrm{Pic}}\nolimits}
\def\Picardstack{\mathcal{P}\!\mathit{ic}}
\def\Picardfunctor{\mathrm{Pic}}
\def\Deformationcategory{\mathcal{D}\!\mathit{ef}}
\usepackage{cleveref}
\crefname{originaldefinition}{original definition}{original definitions}
\Crefname{originaldefinition}{Original definition}{Original definitions}
\makeatletter
\newlabel{definition-etale}{{definition-etale}{source}{}{sourcepointer-1}{}}
\newlabel{definition-etale@cref}{{[originaldefinition][0][]definition-etale}{[1][0][]source}{}{}{}}
\makeatother
\crefname{originallemma}{original lemma}{original lemmas}
\Crefname{originallemma}{Original lemma}{Original lemmas}
\makeatletter
\newlabel{lemma-DM}{{lemma-DM}{source}{}{sourcepointer-2}{}}
\newlabel{lemma-DM@cref}{{[originallemma][0][]lemma-DM}{[1][0][]source}{}{}{}}
\makeatother
\crefname{originallemma}{original lemma}{original lemmas}
\Crefname{originallemma}{Original lemma}{Original lemmas}
\makeatletter
\newlabel{lemma-check-separated-covering}{{lemma-check-separated-covering}{source}{}{sourcepointer-3}{}}
\newlabel{lemma-check-separated-covering@cref}{{[originallemma][0][]lemma-check-separated-covering}{[1][0][]source}{}{}{}}
\makeatother
\crefname{originallemma}{original lemma}{original lemmas}
\Crefname{originallemma}{Original lemma}{Original lemmas}
\makeatletter
\newlabel{lemma-compose-after-separated}{{lemma-compose-after-separated}{source}{}{sourcepointer-4}{}}
\newlabel{lemma-compose-after-separated@cref}{{[originallemma][0][]lemma-compose-after-separated}{[1][0][]source}{}{}{}}
\makeatother
\crefname{originallemma}{original lemma}{original lemmas}
\Crefname{originallemma}{Original lemma}{Original lemmas}
\makeatletter
\newlabel{lemma-etale-smooth-local-source-target}{{lemma-etale-smooth-local-source-target}{source}{}{sourcepointer-5}{}}
\newlabel{lemma-etale-smooth-local-source-target@cref}{{[originallemma][0][]lemma-etale-smooth-local-source-target}{[1][0][]source}{}{}{}}
\makeatother
\crefname{originalreference}{original reference}{original references}
\Crefname{originalreference}{Original reference}{Original references}
\makeatletter
\newlabel{spaces-descent-lemma-descending-property-etale}{{spaces-descent-lemma-descending-property-etale}{source}{}{sourcepointer-6}{}}
\newlabel{spaces-descent-lemma-descending-property-etale@cref}{{[originalreference][0][]spaces-descent-lemma-descending-property-etale}{[1][0][]source}{}{}{}}
\makeatother
\crefname{originaltheorem}{original theorem}{original theorems}
\Crefname{originaltheorem}{Original theorem}{Original theorems}
\makeatletter
\newlabel{theorem-DM}{{theorem-DM}{source}{}{sourcepointer-7}{}}
\newlabel{theorem-DM@cref}{{[originaltheorem][0][]theorem-DM}{[1][0][]source}{}{}{}}
\makeatother
\title{Three equivalent characterizations of etaleness for a morphism of algebraic stacks}
\author{The Stacks Project Authors}
\date{stacks-project:stacks-morphisms.tex:local-snapshot-20261012; GNU FDL 1.2 or later}
\begin{document}
\maketitle
\noindent\textit{Editorial note.} All delivered mathematical text is the author's own, copied from the pinned Stacks Project chapter source sources/stacks-algebraic-stacks-morphisms/pinned-original-source.tex; no statement or proof step was rewritten, reordered or removed.\par\medskip
\noindent\textit{Editorial note.} Preamble provenance: the Stacks Project keeps one single shared preamble for every chapter; this card takes it directly from the staged sources/stacks-algebraic-stacks-morphisms/preamble.tex as three verbatim spans (lines 8-12, 14-18 and 154-217), so no derived file was written and nothing was added or redefined. The wrapper class and the base packages are supplied by the assembler.\par\medskip
\noindent\textit{Editorial note.} Bibliography provenance: neither the counted statement nor its proof contains a citation (this lane rescans the whole span before writing the card), so the card carries no bibliography and no entry was reconstructed.\par\medskip
\noindent\textit{Editorial note.} Selection provenance: this card was opened under the file-based control contract (CONTROL.md), with the mechanical pre-card gate run on its statement and proof interval before the card was written.\par\medskip
\noindent\textit{Allowed standard prerequisites (editorial).}\begin{itemize}
\item The characterization of unramified morphisms through etale covers by algebraic spaces. Invoked inside the counted proof (source lines 8148-8160). Supplies the etale presenting spaces needed for the equivalences.
\item The statement that an unramified flat morphism of finite presentation of algebraic spaces is etale. Invoked inside the counted proof (source lines 8161-8175) as a cross-chapter pointer. Converts the two properties into etaleness for the presenting morphism.
\item The locality of flatness and of local finite presentation in the smooth or fppf topology. Invoked inside the counted proof (source lines 8176-8190). Allows the properties to be moved between the stack and its presentations.
\end{itemize}\medskip
\ifdefined\hypertarget\hypertarget{sourcepointer-1}{}\fi
\noindent\textit{Source reference.} Original definition definition-etale: sources/stacks-algebraic-stacks-morphisms/pinned-original-source.tex, source lines 8078--8078. Navigation pointer only: the counted text refers to the Stacks item definition-etale, which is an existing labeled result of this same chapter and is not part of this card. The pointer carries the target's own label as its number and its definition line as its span; its content is not reproduced here.\par\medskip
\ifdefined\hypertarget\hypertarget{sourcepointer-2}{}\fi
\noindent\textit{Source reference.} Original lemma lemma-DM: sources/stacks-algebraic-stacks-morphisms/pinned-original-source.tex, source lines 4746--4746. Navigation pointer only: the counted text refers to the Stacks item lemma-DM, which is an existing labeled result of this same chapter and is not part of this card. The pointer carries the target's own label as its number and its definition line as its span; its content is not reproduced here.\par\medskip
\ifdefined\hypertarget\hypertarget{sourcepointer-3}{}\fi
\noindent\textit{Source reference.} Original lemma lemma-check-separated-covering: sources/stacks-algebraic-stacks-morphisms/pinned-original-source.tex, source lines 727--727. Navigation pointer only: the counted text refers to the Stacks item lemma-check-separated-covering, which is an existing labeled result of this same chapter and is not part of this card. The pointer carries the target's own label as its number and its definition line as its span; its content is not reproduced here.\par\medskip
\ifdefined\hypertarget\hypertarget{sourcepointer-4}{}\fi
\noindent\textit{Source reference.} Original lemma lemma-compose-after-separated: sources/stacks-algebraic-stacks-morphisms/pinned-original-source.tex, source lines 929--929. Navigation pointer only: the counted text refers to the Stacks item lemma-compose-after-separated, which is an existing labeled result of this same chapter and is not part of this card. The pointer carries the target's own label as its number and its definition line as its span; its content is not reproduced here.\par\medskip
\ifdefined\hypertarget\hypertarget{sourcepointer-5}{}\fi
\noindent\textit{Source reference.} Original lemma lemma-etale-smooth-local-source-target: sources/stacks-algebraic-stacks-morphisms/pinned-original-source.tex, source lines 7765--7765. Navigation pointer only: the counted text refers to the Stacks item lemma-etale-smooth-local-source-target, which is an existing labeled result of this same chapter and is not part of this card. The pointer carries the target's own label as its number and its definition line as its span; its content is not reproduced here.\par\medskip
\ifdefined\hypertarget\hypertarget{sourcepointer-6}{}\fi
\noindent\textit{Source reference.} Original reference spaces-descent-lemma-descending-property-etale: sources/stacks-algebraic-stacks-morphisms/pinned-original-source.tex, source lines 8066--8066. Navigation pointer only: the counted text refers to the Stacks item spaces-descent-lemma-descending-property-etale, which belongs to another chapter of the project and is not part of this card. The pointer carries the target's own label as its number because this chapter's staged snapshot does not contain the other chapter's tag table; its content is not reproduced here.\par\medskip
\ifdefined\hypertarget\hypertarget{sourcepointer-7}{}\fi
\noindent\textit{Source reference.} Original theorem theorem-DM: sources/stacks-algebraic-stacks-morphisms/pinned-original-source.tex, source lines 4589--4589. Navigation pointer only: the counted text refers to the Stacks item theorem-DM, which is an existing labeled result of this same chapter and is not part of this card. The pointer carries the target's own label as its number and its definition line as its span; its content is not reproduced here.\par\medskip
\begin{lemma}
\label{lemma-etale}
Let $f : \mathcal{X} \to \mathcal{Y}$ be a morphism of algebraic stacks.
The following are equivalent
\begin{enumerate}
\item $f$ is \'etale,
\item $f$ is DM and for any morphism $V \to \mathcal{Y}$
where $V$ is an algebraic space and any \'etale morphism
$U \to V \times_\mathcal{Y} \mathcal{X}$ where $U$ is an algebraic space,
the morphism $U \to V$ is \'etale,
\item there exists some surjective, locally of finite presentation, and flat
morphism $W \to \mathcal{Y}$ where $W$ is an algebraic space and some
surjective \'etale morphism $T \to W \times_\mathcal{Y} \mathcal{X}$
where $T$ is an algebraic space such that the morphism $T \to W$ is \'etale.
\end{enumerate}
\end{lemma}

\begin{proof}
Assume (1). Then $f$ is DM and since being \'etale is preserved
by base change, we see that (2) holds.

\medskip\noindent
Assume (2). Choose a scheme $V$ and a surjective \'etale morphism
$V \to \mathcal{Y}$. Choose a scheme $U$ and a surjective \'etale morphism
$U \to V \times_\mathcal{Y} \mathcal{X}$ (Lemma \ref{lemma-DM}).
Thus we see that (3) holds.

\medskip\noindent
Assume $W \to \mathcal{Y}$ and $T \to W \times_\mathcal{Y} \mathcal{X}$
are as in (3). We first check $f$ is DM. Namely, it suffices to check
$W \times_\mathcal{Y} \mathcal{X} \to W$ is DM, see
Lemma \ref{lemma-check-separated-covering}.
By Lemma \ref{lemma-compose-after-separated}
it suffices to check $W \times_\mathcal{Y} \mathcal{X}$ is DM.
This follows from the existence of $T \to W \times_\mathcal{Y} \mathcal{X}$
by (the easy direction of) Theorem \ref{theorem-DM}.

\medskip\noindent
Assume $f$ is DM and $W \to \mathcal{Y}$ and
$T \to W \times_\mathcal{Y} \mathcal{X}$ are as in (3).
Let $V$ be an algebraic space, let $V \to \mathcal{Y}$ be surjective smooth,
let $U$ be an algebraic space, and let
$U \to V \times_\mathcal{Y} \mathcal{X}$ is surjective and \'etale
(Lemma \ref{lemma-DM}). We have to check that $U \to V$ is \'etale.
It suffices to prove $U \times_\mathcal{Y} W \to V \times_\mathcal{Y} W$
is \'etale by Descent on Spaces, Lemma
\ref{spaces-descent-lemma-descending-property-etale}.
We may replace $\mathcal{X}, \mathcal{Y}, W, T, U, V$ by
$\mathcal{X} \times_\mathcal{Y} W, W, W, T, U \times_\mathcal{Y} W,
V \times_\mathcal{Y} W$ (small detail omitted).
Thus we may assume that $Y = \mathcal{Y}$ is an algebraic space, there exists
an algebraic space $T$ and a surjective \'etale morphism
$T \to \mathcal{X}$ such that $T \to Y$ is \'etale, and $U$ and $V$
are as before. In this case we know that
$$
U \to V\text{ is \'etale}
\Leftrightarrow
\mathcal{X} \to Y\text{ is \'etale}
\Leftrightarrow
T \to Y\text{ is \'etale}
$$
by the equivalence of properties (1) and (2) of
Lemma \ref{lemma-etale-smooth-local-source-target}
and Definition \ref{definition-etale}.
This finishes the proof.
\end{proof}

\end{document}
